%==============================================================================
%  CAPÍTULO — CRÉDITOS LABORALES Y PREVISIONALES EN EL CONCURSO
%==============================================================================
\chapter{Créditos laborales y previsionales: verificación, preferencia y pago}
\label{cap:creditos-laborales}

\fichaCapitulo{El primer lugar en la prelación.}{Aplicación concursal del artículo 2472
del Código Civil, topes de las preferencias laborales, pago administrativo, reservas y
tratamiento de las cotizaciones impagas.}

%==============================================================================
\bloqueTeorico
%==============================================================================

\subsection{Por qué el crédito laboral se paga primero}

La preferencia de los créditos laborales no responde a una técnica de garantía, sino a
una decisión de política legislativa: el trabajador no eligió a su deudor, no pudo
diversificar el riesgo ni tarificarlo, y su crédito tiene naturaleza alimentaria. Frente
al acreedor financiero ---que evaluó, tarificó y garantizó su riesgo--- la posición del
trabajador es estructuralmente distinta, y el ordenamiento la reconoce situándolo en la
primera clase del \artcc{2472} \citeyearpar{codigo-civil}.

\subsection{Los límites de la preferencia}

La preferencia no es ilimitada. El legislador la acota mediante topes, y la razón es
igualmente comprensible: una preferencia sin techo permitiría que remuneraciones muy
altas o indemnizaciones convencionales generosas absorbieran íntegramente el activo,
desplazando no sólo a los acreedores valistas sino a los propios trabajadores de menor
renta. Los topes cumplen así una función distributiva dentro de la propia clase.

\subsection{La tensión con las garantías reales}

Los créditos de primera clase se pagan con preferencia sobre los bienes del deudor, pero
la relación con los acreedores hipotecarios y prendarios ---de segunda y tercera clase
sobre bienes determinados--- da lugar a una de las cuestiones más técnicas del derecho
concursal: cómo concurren los créditos de primera clase, que gravan \emph{todo} el
patrimonio, sobre un bien \emph{determinado} afecto a hipoteca o prenda.

La regla del \artcc{2478} fija el orden: los créditos de primera clase no se extienden a
las fincas hipotecadas sino cuando los demás bienes del deudor no alcanzan para pagarlos.
Es decir, \destacar{la primera clase se paga primero con el resto del patrimonio}, y sólo si
éste es insuficiente alcanza al bien hipotecado, con preferencia sobre el acreedor
hipotecario.

\begin{foco}[Por qué esto importa al calcular la recuperación laboral]
El trabajador de primera clase no cobra automáticamente sobre el inmueble hipotecado: cobra
primero sobre los bienes libres. Si éstos bastan, el acreedor hipotecario conserva íntegra su
garantía; si no bastan, el déficit laboral se satisface sobre la finca hipotecada con
prelación. Estimar la recuperación laboral exige, por tanto, mapear antes qué activo está
libre y qué activo está gravado.
\end{foco}

%==============================================================================
\bloqueNormativo
%==============================================================================

\subsection{Categorías de créditos laborales y su tratamiento}

\begin{table}[htbp]
\centering
\small
\caption{Créditos laborales y previsionales en la prelación concursal}
\label{tab:creditos-laborales}
\begin{tabularx}{\textwidth}{@{}>{\raggedright\arraybackslash}p{0.28\textwidth} p{0.16\textwidth} X@{}}
\toprule
\textbf{Crédito} & \textbf{Clase (art. 2472 CC)} & \textbf{Observaciones} \\
\midrule
Remuneraciones devengadas & Primera clase, N.\textsuperscript{o} 5 & Sin tope específico dentro
del numeral \\
\addlinespace
Asignaciones familiares & Primera clase, N.\textsuperscript{o} 5 & — \\
\addlinespace
Cotizaciones previsionales y de salud & Primera clase, N.\textsuperscript{o} 5 & Retenidas y no
enteradas: régimen propio y responsabilidades autónomas \\
\addlinespace
Indemnización del \artcdt{163 bis} \Nro 2 (término concursal) & Primera clase, N.\textsuperscript{o} 5
& \destacar{Tope: 90 UF}; el exceso es valista \\
\addlinespace
Alimentos legales adeudados & Primera clase, N.\textsuperscript{o} 5 & \destacar{Tope: 120 UF}; el
exceso es valista \\
\addlinespace
Feriado proporcional y pendiente & Primera clase & Se devenga hasta la fecha del término \\
\addlinespace
Indemnizaciones legales y convencionales por años de servicio & Primera clase, N.\textsuperscript{o} 8
& \destacar{Tope: 3 ingresos mínimos mensuales por año de servicio y fracción superior a seis
meses, con máximo de 11 años}; el exceso es valista \\
\addlinespace
Créditos laborales controvertidos en juicio & Contingentes & Deben reservarse hasta
sentencia firme \\
\bottomrule
\end{tabularx}
\end{table}

Los topes provienen del texto vigente del \artcc{2472}: el N.\textsuperscript{o} 5 sujeta la
indemnización del \artcdt{163 bis} \Nro 2 a un límite de \destacar{90 UF} y los alimentos legales a
\destacar{120 UF} ---considerándose valista el exceso---; y el N.\textsuperscript{o} 8 limita las
indemnizaciones legales y convencionales de origen laboral a \destacar{tres ingresos mínimos
mensuales por cada año de servicio y fracción superior a seis meses, con un máximo de once años}.
Ambos límites se calculan de forma independiente.

\subsection{Las cotizaciones previsionales impagas}

Las cotizaciones descontadas de la remuneración y no enteradas en las instituciones
previsionales no son, en rigor, patrimonio del empleador: son dineros ajenos retenidos.
Esa naturaleza explica su tratamiento reforzado y las responsabilidades personales que
pueden alcanzar a los administradores.

La \superir{} ha impartido instrucciones específicas sobre pagos y empozamientos ante las
administradoras de fondos de pensiones y la administradora de fondos de cesantía
\parencite{oficio-334, oficio-5003}, y sobre la entrada en vigencia de las modificaciones
al régimen previsional y al seguro de cesantía \parencite{oficio-15795}.

\subsection{El aporte al seguro de cesantía y su imputación}

El aporte del empleador al seguro de cesantía se imputa a la indemnización por años de
servicio. En sede concursal la operación de esta imputación ha requerido precisión
administrativa \parencite{oficio-5080}, por su incidencia directa en el monto que
finalmente debe pagarse con cargo a la masa.

\begin{foco}[Un error de cálculo frecuente]
Omitir la imputación del aporte al seguro de cesantía infla el crédito laboral preferente
y, con ello, reduce indebidamente lo disponible para el resto de la masa. El liquidador
que reparte sobre una base mal calculada responde de la diferencia.
\end{foco}

\subsection{Pago o reserva de créditos laborales}

La procedencia del pago inmediato de los créditos laborales, frente a la alternativa de
constituir reserva, ha sido objeto de pronunciamientos reiterados de la \superir{}
\parencite{oficio-6499, oficio-1765}. La cuestión se plantea cuando el crédito consta
suficientemente pero existen créditos de igual o mejor derecho aún no determinados.

%==============================================================================
\bloquePractico
%==============================================================================

\subsection{Verificación del crédito laboral}

\begin{checklist}[Antecedentes que debe acompañar el trabajador]
  \item Contrato de trabajo y sus anexos.
  \item Liquidaciones de remuneración del período adeudado.
  \item Finiquito, si se hubiere emitido.
  \item Certificado de cotizaciones previsionales.
  \item Sentencia o acta de conciliación, si el crédito consta en juicio.
  \item Instrumento colectivo, si se invocan indemnizaciones convencionales.
  \item Cálculo detallado del crédito, desglosado por concepto y con indicación de la
  preferencia invocada respecto de cada partida.
\end{checklist}

\begin{practica}[Desglose por concepto]
La verificación laboral debe desglosarse partida por partida, porque cada una tiene tope
y preferencia distintos. Una verificación por monto global obliga al liquidador a
reconstruir el cálculo y expone al trabajador a objeciones evitables.
\end{practica}

\subsection{Orden de trabajo del liquidador}

\begin{checklist}[Secuencia de determinación y pago]
  \item Construcción de la nómina laboral con antecedentes documentales.
  \item Cálculo individual por trabajador, desglosado por concepto.
  \item Imputación del aporte al seguro de cesantía a la indemnización por años de
  servicio.
  \item Aplicación de los topes legales a cada partida.
  \item Identificación del exceso sobre los topes, que se reclasifica como crédito
  valista.
  \item Determinación de las contingencias: fueros, juicios en curso, créditos
  controvertidos.
  \item Constitución de reservas por las contingencias identificadas.
  \item Pago administrativo de los créditos suficientemente acreditados.
  \item Empozamiento de los montos no cobrados, conforme a las instrucciones vigentes
  \parencite{oficio-5003}.
\end{checklist}

\subsection{El trabajador que no comparece}

Los créditos de trabajadores que no verifican ni comparecen no desaparecen: deben ser
considerados por el liquidador sobre la base de los registros de la empresa, y los montos
no cobrados quedan sujetos al régimen de empozamiento.

%==============================================================================
\bloqueErrores
%==============================================================================

\errorfrecuente{Verificar el crédito laboral por un monto global}
{El liquidador no puede aplicar los topes por partida y objeta la verificación completa.}
{Desglosar por concepto, con el cálculo y la preferencia invocada respecto de cada uno.}

\errorfrecuente{Omitir la imputación del aporte al seguro de cesantía}
{Se infla el crédito preferente y se perjudica al resto de la masa.}
{Aplicar la imputación conforme a las instrucciones de la \superir{}
\parencite{oficio-5080}.}

\errorfrecuente{Tratar como preferente el exceso de las indemnizaciones convencionales}
{Se altera el reparto en perjuicio de los acreedores valistas y de los propios
trabajadores de menor renta.}
{Reclasificar el exceso sobre el tope como crédito valista.}

\errorfrecuente{Pagar sin reservar los créditos controvertidos}
{Si prosperan las demandas laborales pendientes, no hay fondos con qué pagarlas.}
{Constituir reserva por toda contingencia laboral antes del primer reparto.}

\errorfrecuente{Ignorar las cotizaciones retenidas y no enteradas}
{Subsisten responsabilidades autónomas, incluso personales, de los administradores.}
{Identificarlas desde la construcción de la nómina y darles el tratamiento reforzado que
les corresponde.}
